% !TEX TS-program = lualatex+se+debug
% !TEX encoding = UTF-8 Unicode

\documentclass[14pt, twoside, a4paper]{extreport}
% -----------------------------------------------------------------------------------
% ------------------------- USTAWIENIA STRONY--------------------------
% -----------------------------------------------------------------------------------
\usepackage{geometry}
\geometry{
	left=10mm,
	right=10mm,
	top=15mm,
	bottom=15mm,
	marginparwidth=0mm,
	marginparsep=0mm
	}
\newcommand{\wstawplik}[1]{\input{../include/#1}}%skraca kod
\wstawplik{headerofficium}
% elementy zmienne co do headerofficium:
\graphicspath{{../img/}} % katalog obrazów
%\svgpath{{../include/}}
\grechangestaffsize{20}%rozmiar wyniku GregorioTEX
% koniec 
\gresetlastline{ragged}% justowanie do lewej
% debugowanie:
% %\iffalse %*** %\fi %***
% \tracingall przed problematycznym kodem - bardzo dużo informacji
% lualatex+se+debug - dodatkowe opcje: -file-line-error i -halt-on-error 
% \errorcontextlines
%\def\mydebug{1}
\begin{document}
\makeatletter
\grechangedim{interwordspacetext}{\unexpanded{\the\fontdimen2\font plus \the\gluestretch\gre@space@skip@interwordspacenotes minus \the\glueshrink\gre@space@skip@interwordspacenotes}}{fixed}%
\makeatother
\pagestyle{empty}
%\pagestyle{plain}	
\large
\begin{spacing}{0.7}
\begin{center}
	\section*{Gloria laus}
	%\small\textit{Fontevraud Gradual\\}
	\small\textit{Liber Usualis, Solesmes, 1961, p. 586}
	\vspace{4mm}
\end{center}	
\end{spacing}
\greannotation{Hymn.} 
\greannotation{1} 
\grecommentary[7pt]{Theodulphus, episcopus Aurelianensis † 821}
\gresetinitiallines{1}\gabcsnippet{(c4) <alt>Chorus :</alt>GLó(h)ri(g)a,(gh) laus(gg) et(f) ho(g)nor(gh) ti(f)bi(e) sit,(d.) (,) Rex(g) Chri(g_[oh:h])ste(e) Red(g)ém(gg)ptor :(f.) (:) Cu(fh)i(h) pu(h)e(hvGF')rí(g)le(gh) de(fe)cus(d_f) (,) prom(f)psit(fe) Ho(d)sán(dc~)na(ef) pi(d.)um.(d.) (::Z)
<c><v>\footnotesize{1. }</v></c> Is(h)ra(g)ël(h) es(h!iwj) tu(ji) Rex,(h_g) (,) Da(gh)ví(j)dis(jk) et(j) ín(ji~)cli(h)ta(g) pro(h)les :(h.) (:) Nó(h)mi(h)ne(g) qui(h') in(iwj) Dó(j)mi(ji)ni,(h_g) (,) Rex(gh) be(j)ne(jk)dí(ji)cte,(h) ve(gh)nis.(h.) (::) Gló(h)ri(g)a(gh) (::)
<c><v>\footnotesize{2. }</v></c> Cœ(h)tus(g) in(h) ex(h!iwj)cél(ji)sis(h_g) (,) te(gh) lau(j)dat(jk) caé(ji)li(h)cus(g) o(h)mnis,(h.) (:) Et(h) mor(h)tá(g)lis(h!iwj) ho(ji)mo,(h_g) (,) et(gh) cun(j)cta(jk) cre(j)á(ji)ta(h) si(gh)mul.(h.) (::) Gló(h)ri(g)a(gh) (::)
<c><v>\footnotesize{3. }</v></c> Plebs(h) He(g)braé(h)a(h!iwj) ti(ji)bi(h_g) (,) cum(gh) pal(j)mis(jk) ób(ji~)vi(h)a(g) ve(h)nit :(h.) (:) Cum(h) pre(h)ce,(g) vo(h)to,(h!iwj) hy(ji)mnis,(h_g) (,) ád(gh)su(j)mus(jk) ec(ji)ce(h) ti(gh)bi.(h.) (::) Gló(h)ri(g)a(gh) (::)
<c><v>\footnotesize{4. }</v></c> Hi(h) ti(g)bi(h) pas(h!iwj)sú(ji)ro(h_g) (,) sol(gh)vé(j)bant(jk) mú(ji)ni(h)a(g) lau(h)dis :(h.) (:) Nos(h) ti(h)bi(g) re(h!iwj)gnán(ji)ti(h_g) (,) pán(gh)gi(j)mus(jk) ec(ji)ce(h) me(gh)los.(h.) (::) Gló(h)ri(g)a(gh) (::)
<c><v>\footnotesize{5. }</v></c> Hi(h) pla(h)cu(g)é(h)re(h!iwj) ti(ji)bi,(h_g) (,) plá(gh)ce(j)at(jk) de(j)vó(ji)ti(h)o(g) no(h)stra :(h.) (:) Rex(h) bo(h)ne,(g) Rex(h!iwj) cle(ji)mens,(h_g) (,) cui(gh) bo(j)na(jk) cun(ji~)cta(h) pla(gh)cent.(h.) (::) Gló(h)ri(g)a(gh) (::)}
\vspace{2mm}
\setstretch{1}% odstęp między liniami
\small
\mred{R:}~Chwała, sława, i cześć, niechaj będzie tobie Królu Chryste Odkupicielu: któremu pochwała chłopięca głosiła pobożne Hosanna.
\mred{1.}~Tyś Królem Izraela, i Dawida sławnym potomkiem: Królu błogosławiony, który przychodzisz w Imię Pańskie. 
\mred{2.}~Na wysokościach wysławiają cię wszystkie zastępy niebieskie, i śmiertelny człowiek, oraz wszelkie stworzenia społem.
\mred{3.}~Lud żydowski wyszedł ci naprzeciw z palmami: my stajemy przed tobą z modlitwą, prośbą, pieśniami.
\mred{4.}~Oni tobie idącemu na mękę oddawali powinność chwały: my tobie królującemu śpiewamy oto pieśń.
\mred{5.}~Oni ci się podobali, niechaj (ci) się podoba (także) nasze nabożeństwo: Królu dobry, Królu łaskawy, któremu podoba się wszystko, (co) dobre.
\end{document}